\subsection{应用：Poincaré-Volterra 定理}

定理 1. 设 X 是满足公理 (O \(_{III}\) )（但不一定是豪斯多夫的）的拓扑空间，且设 X 连通且局部连通。设 Y 是拓扑具有可数基的拓扑空间，p: X → Y 是连续映射，使得对每个 y ∈ Y，\(\overline{p}^{1}(y)\) 是 X 的离散子空间。最后设 B 是由 X 的子集组成的一个集，其诸成员的内部覆盖 X，且满足：

\begin{enumerate}
\def\labelenumi{(\roman{enumi})}
\tightlist
\item
  \(p\) 在每个 \(V \in \mathfrak{B}\) 上的限制是 \(V\) 到 \(Y\) 中的一个闭映射。\\
\item
  每个 \(V \in \mathfrak{B}\) 都有一个在 \(V\) 中稠密的可数子集。
\end{enumerate}

则空间 X 是一个可数开集族之并，其中每个开集都含于 B 的某个集合中。

设 \(\mathfrak{B}\) 是 Y 的拓扑的一个可数基。我们称一对 (W, U) 是特异的，如果 (i) \(U \in \mathfrak{B}\)，且 (ii) W 是 \(\overline{p}(U)\) 的一个分支，它含于 \(\mathfrak{B}\) 的某个集合中。

引理 1. 若 \(x\) 是 \(\mathbf{X}\) 的任一点，则存在特异对 \((\mathbf{W}, \mathbf{U})\) 使得 \(x \in \mathbf{W}\)。

逆像 \(\overline{p} (p(x))\) 是离散的，因而存在 \(x\) 在 \(\mathbf{X}\) 中的一个邻域，使得对其每一点 \(x^{\prime}\)（异于 \(x\) 者）都有像 \(p(x^{\prime})\neq p(x)\)；由于 \(\mathbf{X}\) 满足 \((\mathrm{O}_{\mathrm{III}})\)，存在具有该性质的闭邻域 \(\mathbf{V}\)（即 \(x\) 的闭邻域），且还可假设 \(\mathbf{V}\) 含于 \(\mathfrak{B}\) 的某个集合中。设 \(\mathbf{F}\) 是 \(\mathbf{V}\) 在 \(\mathbf{X}\) 中的边界。由定理的条件 (i)，\(p(\mathbf{F})\) 在 \(\mathbf{Y}\) 中是闭的；且由于 \(p(\mathbf{F})\) 不含 \(p(x)\)，存在集合 \(\mathbf{U}\in \mathfrak{B}\)，它含有 \(p(x)\) 且不与 \(p(\mathbf{F})\) 相交。设 \(\mathbf{W}\) 是 \(x\) 在 \(\overline{p}^{1}(\mathbf{U})\) 中的分支；则只须证明 \(\mathbf{W}\subset \mathring{\mathbf{V}}\)。若不然，\(\mathbf{W}\) 将与 \(\mathbf{F}\) 相交（第 1 目，命题 3），因而 \(p(\mathbf{F})\) 将与 \(\mathbf{U}\) 相交，这与 \(\mathbf{U}\) 的定义相矛盾。

引理 2. 若 (W, U) 是一个特异对，则使 W' 与 W 相交的所有特异对 (W', U') 所成的集是可数的。

既然 \(\mathfrak{B}\) 是可数的，只须证明：对给定的 \(\mathbf{U}' \in \mathfrak{B}\)，特异对 \((\mathbf{W}', \mathbf{U}')\)（其中 \(\mathbf{W}'\) 与 \(\mathbf{W}\) 相交者）所成的集是可数的。现在这些集合 \(\mathbf{W}'\) 都是开的，因为 \(\mathbf{X}\) 局部连通（第 6 目，命题 11），且两两不相交，因为它们是 \(\overline{\boldsymbol{p}}(\mathbf{U}')\) 的分支；于是诸集合 \(\mathbf{W}' \cap \mathbf{W}\) 是开的且两两不相交。但 \(\mathbf{W}\) 含有一个在 \(\mathbf{W}\) 中稠密的可数子集；故诸 \(\mathbf{W}'\)（其中 \(\mathbf{W}' \cap \mathbf{W}\) 非空者）所成的集也是可数的。

为证定理 1，考虑 X 的两点 \(x\), \(x'\) 之间的如下关系 R：“存在特异对的一个有限序列 \((W_i, U_i)\)（\(1 \leqslant i \leqslant n\)），使得 \(x \in W_1\)，\(x' \in W_n\)，且 \(W_i \cap W_{i+1} \neq \emptyset\) 对 \(1 \leqslant i \leqslant n - 1\) 成立。”

引理 1 表明 R 是自反的，且容易验证 R 是对称且传递的，故 R 是一个等价关系；又由于诸 \(W_{i}\) 是开的，每个模 R 等价类都在 X 中是开的。但 X 连通；故只能有一个等价类，即 X 的任意两点都模 R 同余。我们将由此推出，X 是特异对第一元素所成的一个可数族之并，这就证明了定理 1。为此，设 x 是 X 的任一点，并对 n 用归纳法定义 X 的开子集的序列 \((\mathbf{C}_{n})\) 如下：由引理 1，存在特异对 \((\mathbf{W}_{1}, \mathbf{U}_{1})\) 使得 \(x \in W_{1}\)，我们取 \(C_{1} = W_{1}\)；若 n \textgreater{} 1，则 \(C_{n}\) 取为所有使 W 与 \(C_{n-1}\) 相交的特异对 (W, U) 的第一元素 W 的并。由引理 2，对 n 用归纳法立即可见，\(C_{n}\) 是特异对第一元素的一个可数并。最后，每个 \(x' \in X\) 都属于某个 \(C_{n}\)；因为存在有限序列

\[
\left(\mathrm{W} _ {i} ^ {\prime}, \mathrm{U} _ {i} ^ {\prime}\right) _ {1 \leqslant i \leqslant m}
\]

的特异对，使得 \(x \in W_1'\)，\(x' \in W_m'\)，且 \(W_i' \cap W_{i+1}' \neq \emptyset\) 对 \(1 \leqslant i \leqslant m - 1\) 成立；对 \(i\) 用归纳法可以看出，\(W_i' \subset C_{i+1}\) 对所有 \(i\) 都成立，从而 \(x' \in C_{m+1}\)。

证毕。

推论 1. 设 Y 是拓扑具有可数基的正则空间（*）。设 X 是连通且局部连通的空间，且设 \(p: \mathbf{X} \to \mathbf{Y}\) 是具有以下性质的连续映射：对每个 \(x \in \mathbf{X}\)，存在 \(x\) 在 X 中的一个闭邻域 V，使得 \(p\) 在 V 上的限制是 V 到 Y 的一个闭子空间上的同胚。则 X 是正则的且 X 的拓扑具有可数基。

首先，诸假设蕴含 X 是正则的（§ 8，第 4 目，命题 13）。让我们证明，若取 B 为 X 的所有这样的闭子集 V 所成的集：p 在 V 上的限制是 V 到 Y 的一个闭子空间上的同胚，则定理 1 的条件得到满足。按假设，B 中诸集合的内部覆盖 Y，且凭关于 Y 的假设，每个 \(V \in B\) 都具有可数基，因而含有一个可数稠密子集（§ 1，第 6 目，命题 6）。此外，若 \(x \in \overline{p}(y)\)，则存在 x 在 X 中的一个邻域 \(V \in B\)，使得 V 不含 \(\overline{p}(y)\) 的异于 \(x\) 的任何点，从而 \(\overline{p}(y)\) 是一个离散空间。于是可以应用定理 1，它表明 X 是开集的一个可数族 \((\mathrm{T}_n)_{n\geqslant 0}\) 之并，使得每个子空间 \(\mathrm{T}_n\) 都具有可数基 \((\mathrm{U}_{mn})_{m\geqslant 0}\)。于是诸 \(\mathrm{U}_{mn}\)（\(m\geqslant 0\)，\(n\geqslant 0\)）构成 X 的拓扑的一个基（§ 3，第 1 目，注）。

推论 2. 设 X 局部紧、连通且局部连通，并设 X 的每一点都有一个具有可数基的邻域。设 Y 是拓扑具有可数基的豪斯多夫空间，且设 \(p: \mathbf{X} \to \mathbf{Y}\) 是连续映射，使得对每个 \(y \in \mathbf{Y}\)，\(\overline{p}^1(y)\) 是 X 的离散子空间。则 X 的拓扑具有可数基。

对每个 \(x \in X\)，设 \(V_{x}\) 是 x 在 X 中的一个具有可数基的紧邻域。由 §9，第 4 目，定理 2，推论 2 可知，诸 \(V_{x}\) 所成的集 B 满足定理 1 的条件，于是我们像推论 1 中那样完成证明。

注意，在此推论中，可能发生这样的情形：p 在 X 的某一点的任意小邻域 V 上的限制不是 V 到 \(p(\mathrm{V})\) 上的同胚。

推论 3（Poincaré-Volterra 定理）. 设 Y 是局部紧、局部连通且拓扑具有可数基的空间。设 X 是连通豪斯多夫空间，且设 \(p: \mathbf{X} \to \mathbf{Y}\) 是具有以下性质的连续映射：对每个 \(x \in \mathbf{X}\)，存在 \(x\) 在 X 中的一个开邻域 U，使得 \(p\) 在 U 上的限制是 U 到 Y 的一个开子空间上的同胚。则 X 局部紧且局部连通，且 X 的拓扑具有可数基。

显然 X 局部连通。又每个 \(x \in X\) 都有 X 中的一个开邻域 U，使得 p 在 U 上的限制把 U 同胚地映到 Y 的开子空间 \(p(U)\) 上。由于 \(p(U)\) 是 Y 的局部紧子空间（\(§ 9, no. 7, Proposition 13\)），存在 \(p(x)\) 的一个紧邻域 W，它含于 \(p(U)\) 中，从而 \(U \cap \overline{p}(W)\) 是含于 U 中的 x 的一个紧邻域；由于按假设 X 是豪斯多夫的，故 X 局部紧。\(U \cap \overline{p}(W)\) 是紧的，故在 X 中是闭的（\(§ 9, no. 3, Proposition 4\)），因而推论 1 的条件得到满足；故 X 的拓扑具有可数基。

